\section{Discussion}

A facial module is scientifically attractive because it could provide continuous, non-contact information during interaction. The same property creates a high risk of overclaiming. Facial dynamics are affected by task content, speech, gaze direction, lighting, device pose, idiosyncratic expressiveness, and many other factors. Treating them as a direct readout of emotion would make the adaptive system difficult to validate and ethically harder to justify.

The incremental-validity design avoids that mistake. It asks whether facial behavior contributes information after simpler and less intrusive signals are already known. This ordering is important because a technically impressive predictor can still be unnecessary if task behavior and direct child report perform equally well.

The main limitation is label validity. Challenge level, self-reported difficulty, recovery status, and stop requests are not interchangeable constructs. Each target therefore requires a separate analysis and should not be collapsed into a latent true affect variable without independent validation.

A second limitation is ecological robustness. Home use introduces uncontrolled lighting, camera angle, motion, and partial occlusion. Missingness and tracking failure must be modeled as part of feasibility rather than excluded until the dataset becomes artificially clean.

The strongest result for this paper would not necessarily be a highly accurate facial model. A well-powered null result showing that facial sensing adds little beyond simpler telemetry would provide a valuable design conclusion and reduce unnecessary sensing in child-facing digital health systems.
